\newcommand{\ModelSentence}{\texttt{alibaba/qwen3-coder-30b-a3b} (Qwen3-Coder-30B-A3B-Instruct, an open-weight coding model with a fixed release, accessed through the Vercel AI Gateway)}
\newcommand{\RunHardware}{Agent runs for Python projects were executed on Google Colab (Python 3.13); every scoring run builds a fresh virtual environment. C++ and Java projects are scored in DevBench's container image in GitHub Codespaces.}
\newcommand{\AgentAbstract}{In a pre-registered pilot on 4 Python projects, team designs did not differ in correctness, but where the criterion recommends one agent (3 projects), a single agent was as correct as every split team and used 34--90\% fewer tokens.}
\newcommand{\AgentConclusion}{A pre-registered pilot on 4 Python projects found no correctness difference between team designs, so its go/no-go rules return a provisional NO-GO; but on all 3 projects for which the criterion recommends one agent, following that recommendation was as correct as every split team and cut token use by 34--90\%. The main study will use tasks on which models pass a meaningful share of tests, add a second model, and test H1--H4.}
\newcommand{\ContributionFour}{A pre-registered agent study with four hypotheses, controlled conditions and an equal token budget per run, a zero-cost verification of the agent harness, and a pilot reported in full, including its go/no-go outcome (Section~\ref{sec:agents}).}
\newcommand{\ThreatExtra}{DevBench's tests also check identifiers that are absent from the requirements and the design, which lowers every condition's absolute score; TextCNN is excluded from the pilot because its reference code passes only 5 of 13 tests in the sandbox, as the pre-registered scorer rule requires.}
\newcommand{\PilotDiscussion}{\paragraph{What the pilot shows.} Projects for which the criterion recommends a single agent are where it differs most from graph partitioners. Of the 8 projects with more than one class for which it recommends a single agent, one agent per class splits all 8, Louvain splits 7 and Infomap, the partitioner used by Co-Coder, splits 2 (geotext and particle-swarm-optimization), in both cases into the same teams as one agent per class. Experiment 1b runs exactly these 2 projects, so the single agent is compared with the teams that Infomap, one-agent-per-class and random splitting would use. The pilot was meant to expose problems before the main study, and it did. Its harness checks found and fixed four defects that would otherwise have biased the comparison, and its floor effect shows that, with an open-weight model of this size, DevBench's Python tests are too strict to separate the conditions, because they reward identifiers that the design does not contain. Three changes are therefore proposed for the main study and will be registered before any further run: tasks whose tests call only names that appear in the design, or designs completed with the public interface that the tests call, identically for every condition; a second, stronger model; and token cost and inter-agent messages reported as outcomes alongside tests passed, since the pilot shows that they differ even when correctness does not.}
